% Lösungen Einheit 2 von 5 – Kürzen und Erweitern (zusammenbau v0.1)
\einheitenkopf{Einheit 2 von 5 · Kürzen und Erweitern}

%% TODO Lösung der Erklärzeile 14g) fehlt in der Bank
\erg{14}{a) mit $3$ \quad b) $\frac{12}{20}$ \quad c) $\frac{3}{18}$ \quad d) $\frac{8}{18}$ \quad e) $\frac{15}{21}$ \quad f) $\frac{10}{15}$ \quad g) \ldots \quad h) $\frac{2}{3}$ \quad i) $\frac{16}{24} = \frac{2}{3}$ \quad j) $\frac{15}{20}$ \quad k) $\frac{3}{9}$ und $\frac{5}{9}$ \quad l) $\frac{15}{20}$ und $\frac{8}{20}$ \quad m) $\frac{13}{5} = 2\frac{3}{5}$ \quad n) $\frac{12}{16} = \frac{3}{4} = 75\,\%$}
\erg{15}{a) $4$ Teile, also $\frac{4}{8}$}

15a)

\bruchrechteck{1}{2} \\ \bruchrechteck{4}{8}

\erg{16}{a) $\frac{7}{9}$}
\erg{17}{a) Sie hat nur den Zähler mal $3$ genommen. Richtig: $\frac{2}{7} = \frac{6}{21}$. \quad b) Nein: Es gibt dreimal so viele Teile, jedes ist aber nur ein Drittel so groß – der Anteil bleibt gleich. \quad c) $\frac{6}{15}$ \quad d) $\frac{45}{60} = \frac{3}{4}$ Stunde}
